\documentclass[11pt,a4paper]{article}

% ── Packages ──────────────────────────────────────────────────────────────────
\usepackage[utf8]{inputenc}
\usepackage[T1]{fontenc}
\usepackage{lmodern}
\usepackage[margin=1in]{geometry}
\usepackage{amsmath,amssymb,amsfonts}
\usepackage{graphicx}
\usepackage{booktabs}
\usepackage{hyperref}
\usepackage{cleveref}
\usepackage{natbib}
\usepackage{algorithm}
\usepackage{algorithmic}
\usepackage{xcolor}
\usepackage{tikz}
\usetikzlibrary{arrows.meta,positioning,shapes.geometric}
\usepackage{float}
\usepackage{subcaption}
\usepackage{enumitem}
\usepackage{microtype}
\usepackage{fancyhdr}
\usepackage{abstract}
\usepackage{multirow}

% ── Page Style ────────────────────────────────────────────────────────────────
\pagestyle{fancy}
\fancyhf{}
\fancyhead[L]{\small Zoo Labs Foundation}
\fancyhead[R]{\small Few-Shot Visual Recognition}
\fancyfoot[C]{\thepage}
\renewcommand{\headrulewidth}{0.4pt}

% ── Metadata ──────────────────────────────────────────────────────────────────
\title{%
  \textbf{Few-Shot Visual Recognition for\\Biodiversity Monitoring}\\[0.5em]
  \large Zoo Labs Foundation Research Paper ZOO-CLS-2026
}

\author{
  Zoo Labs Foundation Research\\
  \texttt{research@zoo.ngo}\\[0.5em]
  \small With contributions from the Zoo Open Science Network
}

\date{February 2026}

% ── Macros ────────────────────────────────────────────────────────────────────
\newcommand{\R}{\mathbb{R}}
\newcommand{\E}{\mathbb{E}}
\newcommand{\Loss}{\mathcal{L}}
\newcommand{\Support}{\mathcal{S}}
\newcommand{\Query}{\mathcal{Q}}

\begin{document}

\maketitle

% ── Abstract ──────────────────────────────────────────────────────────────────
\begin{abstract}
\noindent
Automated species identification is fundamental to scalable biodiversity
monitoring, yet most deep learning classifiers require hundreds or
thousands of labeled images per species---a prohibitive requirement for
the vast majority of Earth's estimated 8.7 million species. We present
\textbf{TaxoNet}, a few-shot visual recognition framework designed for
fine-grained species classification in ecological monitoring contexts.
TaxoNet combines a vision transformer backbone pre-trained on large-scale
ecological imagery with a taxonomy-aware meta-learning strategy that
exploits the hierarchical structure of biological classification (kingdom
through species). A novel \emph{taxonomic prototypical network}
constructs class prototypes at multiple taxonomic ranks and fuses them
through learned rank-attention weights, enabling robust identification
from as few as 1--5 reference images per species. We evaluate TaxoNet on
four biodiversity datasets spanning birds, insects, plants, and marine
organisms, achieving state-of-the-art 1-shot accuracy of 78.4\% and
5-shot accuracy of 91.2\% on held-out species. We further demonstrate
integration with citizen science platforms, where TaxoNet assists
non-expert observers in identifying species with 94.7\% accuracy when
combined with geographic and phenological priors. All models and datasets
are released through the Zoo Open Science Network.

\medskip
\noindent\textbf{Keywords:} few-shot learning, fine-grained
classification, biodiversity monitoring, taxonomic hierarchy, citizen
science, species identification
\end{abstract}

\newpage
\tableofcontents
\newpage

% ══════════════════════════════════════════════════════════════════════════════
\section{Introduction}
\label{sec:introduction}
% ══════════════════════════════════════════════════════════════════════════════

Biodiversity monitoring at global scales requires automated species
identification systems capable of operating across the full breadth of
taxonomic diversity \citep{christin2019applications}. Camera traps,
acoustic sensors, citizen science platforms, and environmental DNA (eDNA)
sampling generate petabytes of observational data annually, far exceeding
manual identification capacity \citep{steenweg2017scaling}.

Deep learning has revolutionized image-based species identification, with
convolutional neural networks achieving expert-level performance on
well-studied taxa such as birds \citep{van2015building}, plants
\citep{goeau2022overview}, and mammals \citep{norouzzadeh2018automatically}.
However, these successes depend on large labeled training datasets---a
critical limitation given that:

\begin{enumerate}[label=(\roman*)]
  \item The majority of species have fewer than 50 available images in
    public databases \citep{lorieul2022overview}.
  \item Newly discovered species (approximately 18{,}000 per year
    \citep{costello2013can}) have no training data at all.
  \item Rare and cryptic species---often those of greatest conservation
    concern---are systematically underrepresented in image databases.
  \item Intraspecific visual variation (sexual dimorphism, juvenile
    plumage, seasonal morphs) multiplies the effective number of
    required examples.
\end{enumerate}

Few-shot learning addresses this data scarcity by training models to
recognize novel classes from very few examples
\citep{wang2020generalizing}. While few-shot methods have advanced
rapidly in computer vision, their application to ecological species
identification introduces unique challenges:

\begin{itemize}
  \item \textbf{Fine-grained discrimination}: Species within the same
    genus can differ by subtle visual features (wing venation, petal
    shape, scale patterns) that require high-resolution feature
    extraction.
  \item \textbf{High inter-class similarity}: Closely related species
    form dense clusters in feature space, making prototype-based
    methods fragile.
  \item \textbf{Domain shift}: Training images from curated databases
    differ substantially from field photos taken by citizen scientists
    under variable conditions.
  \item \textbf{Hierarchical structure}: Biological taxonomy provides a
    natural hierarchy that standard few-shot methods ignore.
\end{itemize}

\subsection{Contributions}

We introduce \textbf{TaxoNet}, a few-shot species recognition system
with the following contributions:

\begin{enumerate}
  \item A \textbf{taxonomy-aware prototypical network} that constructs
    multi-rank prototypes (family, genus, species) and fuses them
    through learned attention, enabling few-shot identification to
    leverage phylogenetic structure.

  \item A \textbf{hierarchical meta-learning} training strategy with
    taxonomic episode sampling that ensures the model learns
    discriminative features at every rank.

  \item A \textbf{domain-adaptive pre-training} pipeline using 28
    million ecological images from iNaturalist, GBIF, and museum
    collections.

  \item Integration with \textbf{geographic and phenological priors}
    that improve practical field identification accuracy by
    constraining the candidate species set.

  \item Comprehensive evaluation across \textbf{four major taxonomic
    groups} (birds, insects, plants, marine organisms) with separate
    validation on \textbf{citizen science field data}.
\end{enumerate}

\subsection{Zoo Labs Foundation Context}

This work supports Zoo Labs Foundation's mission of open, decentralized
science infrastructure. TaxoNet is designed for deployment on the Zoo
citizen science platform and integrates with the Foundation's
decentralized data pipelines. Model weights and training code are
released under the Apache 2.0 license, and the evaluation benchmark
is maintained as a community resource through the Zoo Open Science
Network at \texttt{zips.zoo.ngo}.


% ══════════════════════════════════════════════════════════════════════════════
\section{Related Work}
\label{sec:related}
% ══════════════════════════════════════════════════════════════════════════════

\subsection{Species Identification with Deep Learning}

The past decade has seen dramatic advances in automated species
identification. iNaturalist's computer vision system classifies over
76{,}000 species \citep{van2021benchmarking}. eBird's Merlin app
achieves high accuracy for North American birds using audio and image
inputs \citep{kahl2021birdnet}. PlantNet identifies over 38{,}000 plant
species from photographs \citep{goeau2022overview}. These systems rely
on extensive labeled datasets, with performance degrading sharply for
species with fewer than 100 training images.

\subsection{Few-Shot Learning}

Few-shot learning aims to classify novel classes from $K$ labeled
examples per class (a ``$K$-shot'' setting). Metric-based approaches
learn an embedding space where classification reduces to nearest-neighbor
comparison: matching networks \citep{vinyals2016matching}, prototypical
networks \citep{snell2017prototypical}, and relation networks
\citep{sung2018learning}. Optimization-based methods learn
initialization points for rapid adaptation: MAML \citep{finn2017model}
and its variants. More recently, large pre-trained models with
in-context learning have shown strong few-shot capabilities
\citep{brown2020language}.

\subsection{Hierarchical Classification}

Taxonomic hierarchy has been exploited in traditional species
classification through hierarchical softmax \citep{deng2014large},
label embedding with taxonomy trees \citep{akata2015label}, and
curriculum learning from coarse to fine \citep{bengio2009curriculum}.
\citet{chang2021your} incorporate hierarchical labels into contrastive
learning. However, hierarchical few-shot learning remains
underexplored, with only \citet{li2019large} and
\citet{liu2020many} addressing it in non-ecological contexts.

\subsection{Fine-Grained Recognition}

Fine-grained visual classification distinguishes subordinate categories
(species, cultivars, model variants). Key approaches include bilinear
pooling \citep{lin2015bilinear}, part-based attention
\citep{zheng2017learning}, and feature recalibration
\citep{hu2018squeeze}. Vision transformers (ViTs) have become dominant
due to their ability to capture global context
\citep{dosovitskiy2021image, he2022transfg}. Few-shot fine-grained
recognition remains challenging due to the high intra-class and low
inter-class variance \citep{wei2019piecewise}.


% ══════════════════════════════════════════════════════════════════════════════
\section{Problem Formulation}
\label{sec:problem}
% ══════════════════════════════════════════════════════════════════════════════

\subsection{Task Definition}

We consider the standard $N$-way $K$-shot classification problem. Given a
\emph{support set} $\Support = \{(x_i, y_i)\}_{i=1}^{N \times K}$
containing $K$ labeled images for each of $N$ novel species, and a
\emph{query set} $\Query = \{x_j\}_{j=1}^{Q}$ of unlabeled images, the
goal is to classify each query into one of the $N$ species.

\subsection{Taxonomic Structure}

Each species $s$ belongs to a taxonomic hierarchy:
\begin{equation}
  \text{Kingdom} \supset \text{Phylum} \supset \text{Class} \supset
  \text{Order} \supset \text{Family} \supset \text{Genus} \supset
  \text{Species}
\end{equation}

We denote the taxonomic label of species $s$ at rank $r$ as $t_r(s)$.
Two species in the same genus share $t_r$ for $r \in
\{\text{kingdom}, \ldots, \text{genus}\}$ but differ at species rank.

\subsection{Incorporating Priors}

For practical deployment, geographic coordinates $(\text{lat},
\text{lon})$ and observation date $d$ provide strong priors on which
species are plausible at a given location and time. We define the
candidate set:
\begin{equation}
  \mathcal{C}(\text{lat}, \text{lon}, d) = \{s : s \text{ has recorded
    occurrence within } r \text{ km and } \pm \Delta d \text{ days}\}
\end{equation}
which restricts the $N$-way classification to geographically and
temporally plausible species.


% ══════════════════════════════════════════════════════════════════════════════
\section{TaxoNet Architecture}
\label{sec:architecture}
% ══════════════════════════════════════════════════════════════════════════════

\subsection{Overview}

TaxoNet consists of four components: (1) a vision transformer feature
extractor, (2) a taxonomic prototype constructor, (3) a rank-attention
fusion module, and (4) optional geographic/phenological prior
integration.

\subsection{Feature Extractor}

We use a Vision Transformer (ViT-L/14) pre-trained through our
domain-adaptive pipeline (Section~\ref{sec:pretraining}). Input images
are resized to $518 \times 518$ pixels and split into $14 \times 14$
patches. The transformer produces a sequence of patch embeddings, from
which we extract both the \texttt{[CLS]} token embedding $\mathbf{h}
\in \R^{1024}$ and the full patch embedding matrix $\mathbf{H} \in
\R^{1369 \times 1024}$.

For fine-grained discrimination, we apply a learned spatial attention
module that identifies discriminative image regions:
\begin{equation}
  \mathbf{a} = \text{softmax}(\mathbf{W}_a \mathbf{H}^\top) \in
  \R^{M \times 1369}
\end{equation}
where $M = 4$ attention heads focus on different discriminative parts.
The attended feature is:
\begin{equation}
  \mathbf{f} = [\mathbf{h} \| \text{vec}(\mathbf{a} \mathbf{H})]
  \in \R^{1024 + 4 \times 1024} = \R^{5120}
\end{equation}
projected to a 512-dimensional embedding via a linear layer.

\subsection{Taxonomic Prototypical Network}

Given the support set $\Support$ for an $N$-way $K$-shot episode, we
construct prototypes at three taxonomic ranks: family ($r_F$), genus
($r_G$), and species ($r_S$).

\paragraph{Species prototype.} The standard prototypical network
prototype:
\begin{equation}
  \mathbf{c}_{r_S}^{(s)} = \frac{1}{K} \sum_{(x_i, y_i) \in
  \Support : y_i = s} \phi(x_i)
  \label{eq:species_proto}
\end{equation}
where $\phi(\cdot)$ is the feature extractor.

\paragraph{Genus prototype.} Aggregating all species in the same genus:
\begin{equation}
  \mathbf{c}_{r_G}^{(g)} = \frac{1}{|\{s : t_{r_G}(s) = g\}|}
  \sum_{s : t_{r_G}(s) = g} \mathbf{c}_{r_S}^{(s)}
\end{equation}

\paragraph{Family prototype.} Similarly aggregating across genera in
the same family:
\begin{equation}
  \mathbf{c}_{r_F}^{(f)} = \frac{1}{|\{g : t_{r_F}(g) = f\}|}
  \sum_{g : t_{r_F}(g) = f} \mathbf{c}_{r_G}^{(g)}
\end{equation}

\subsection{Rank-Attention Fusion}

For a query embedding $\mathbf{q} = \phi(x)$, we compute distances to
each prototype at each rank. The rank-attention module learns to weight
these distances adaptively:

\begin{equation}
  d_r(x, s) = \|\mathbf{q} - \mathbf{c}_r^{(t_r(s))}\|^2
\end{equation}

The attention weights depend on the query itself:
\begin{equation}
  \boldsymbol{\alpha}(\mathbf{q}) = \text{softmax}(\mathbf{W}_r
  \mathbf{q}) \in \R^3
\end{equation}

The fused distance is:
\begin{equation}
  d_{\text{fused}}(x, s) = \sum_{r \in \{r_F, r_G, r_S\}}
  \alpha_r(\mathbf{q}) \cdot d_r(x, s)
\end{equation}

The predicted class probability is:
\begin{equation}
  p(y = s \mid x) = \frac{\exp(-d_{\text{fused}}(x, s))}
  {\sum_{s'} \exp(-d_{\text{fused}}(x, s'))}
  \label{eq:prediction}
\end{equation}

\subsection{Geographic and Phenological Priors}

When location and date metadata are available, we compute a prior
probability $p_{\text{geo}}(s \mid \text{lat}, \text{lon}, d)$ from
species range maps and phenological data. This prior is combined with
the visual prediction via:
\begin{equation}
  p_{\text{final}}(s \mid x, \text{lat}, \text{lon}, d) \propto
  p(y = s \mid x)^{1-\beta} \cdot
  p_{\text{geo}}(s \mid \text{lat}, \text{lon}, d)^{\beta}
\end{equation}
where $\beta = 0.3$ balances visual and geographic evidence. This
effectively re-ranks predictions by geographic plausibility.


% ══════════════════════════════════════════════════════════════════════════════
\section{Domain-Adaptive Pre-Training}
\label{sec:pretraining}
% ══════════════════════════════════════════════════════════════════════════════

Standard ImageNet pre-training provides suboptimal features for
fine-grained ecological recognition due to domain mismatch (controlled
studio photos vs.\ field conditions) and category mismatch (1{,}000
general categories vs.\ thousands of visually similar species).

\subsection{Pre-Training Data}

We assemble \textbf{EcoVision-28M}, a dataset of 28 million ecological
images from:

\begin{itemize}
  \item \textbf{iNaturalist 2021}: 10 million images across 10{,}000
    species \citep{van2021benchmarking}.
  \item \textbf{GBIF media}: 8 million specimen and field images.
  \item \textbf{Museum collections}: 5 million digitized specimen
    photographs from 14 natural history museums.
  \item \textbf{Camera trap archives}: 3 million images from Wildlife
    Insights and LILA BC \citep{beery2021iwildcam}.
  \item \textbf{Citizen science platforms}: 2 million images from
    eBird, BugGuide, and Reef Life Survey.
\end{itemize}

\subsection{Self-Supervised Pre-Training}

We pre-train ViT-L/14 using DINOv2 \citep{oquab2024dinov2} on
EcoVision-28M for 300{,}000 iterations with batch size 1{,}024. The
self-supervised objective learns representations that capture
fine-grained visual structure without requiring species labels.

\subsection{Supervised Fine-Tuning}

After self-supervised pre-training, we fine-tune on the labeled subset
(iNaturalist 2021, 10M images, 10{,}000 species) using a hierarchical
classification head with losses at family, genus, and species level:
\begin{equation}
  \Loss_{\text{pretrain}} = \Loss_{\text{species}} + 0.3
  \Loss_{\text{genus}} + 0.1 \Loss_{\text{family}}
\end{equation}


% ══════════════════════════════════════════════════════════════════════════════
\section{Hierarchical Meta-Learning}
\label{sec:metalearning}
% ══════════════════════════════════════════════════════════════════════════════

\subsection{Episode Construction}

Standard meta-learning samples episodes uniformly. We introduce
\textbf{taxonomic episode sampling} that ensures episodes contain
discriminative challenges at each rank.

For each training episode, we:
\begin{enumerate}
  \item Sample a focal family $f$ with probability proportional to its
    species richness.
  \item From $f$, sample 2--4 genera.
  \item From each genus, sample 2--5 species.
  \item Include 1--2 ``distractor'' species from the same order but
    different family.
\end{enumerate}

This ensures the model must discriminate at multiple taxonomic levels
within each episode, unlike uniform sampling which often produces
episodes of distantly related, trivially distinguishable species.

\subsection{Multi-Rank Loss}

The training loss combines classification losses at each rank:
\begin{equation}
  \Loss_{\text{meta}} = \Loss_{\text{CE}}^{(r_S)} + \gamma_G
  \Loss_{\text{CE}}^{(r_G)} + \gamma_F \Loss_{\text{CE}}^{(r_F)}
  + \lambda \Loss_{\text{triplet}}
\end{equation}
where $\Loss_{\text{CE}}^{(r)}$ is the cross-entropy loss for
classification at rank $r$, $\Loss_{\text{triplet}}$ is a taxonomic
triplet loss that pushes apart embeddings of different genera while
pulling together congeneric species, and $\gamma_G = 0.3$, $\gamma_F =
0.1$, $\lambda = 0.2$ are weighting coefficients.

\subsection{Taxonomic Triplet Loss}

We define:
\begin{equation}
  \Loss_{\text{triplet}} = \sum_{(a, p, n)} \max(0,
  \|\phi(x_a) - \phi(x_p)\| - \|\phi(x_a) - \phi(x_n)\| + m(a, n))
\end{equation}
where $(a, p, n)$ is an anchor-positive-negative triplet, and the margin
$m(a, n)$ scales with taxonomic distance:
\begin{equation}
  m(a, n) = m_0 \cdot (1 - \text{sim}_{\text{tax}}(a, n))
\end{equation}
Here $\text{sim}_{\text{tax}}(a, n) \in [0, 1]$ is the fraction of
shared taxonomic ranks, and $m_0 = 0.5$ is the base margin.

\subsection{Training Protocol}

\begin{itemize}
  \item \textbf{Meta-training set}: 8{,}000 species (disjoint from
    evaluation species).
  \item \textbf{Meta-validation set}: 1{,}000 species.
  \item \textbf{Episodes per epoch}: 10{,}000.
  \item \textbf{Optimizer}: AdamW, learning rate $5 \times 10^{-5}$,
    cosine schedule over 100 epochs.
  \item \textbf{Episode configuration}: 10-way $K$-shot, $K \in
    \{1, 5\}$, with 15 queries per class.
\end{itemize}


% ══════════════════════════════════════════════════════════════════════════════
\section{Evaluation}
\label{sec:evaluation}
% ══════════════════════════════════════════════════════════════════════════════

\subsection{Datasets}

We evaluate on four domain-specific benchmarks:

\begin{table}[H]
\centering
\caption{Evaluation datasets. ``Novel'' species have no overlap with
training data.}
\label{tab:datasets}
\begin{tabular}{lcccc}
\toprule
\textbf{Dataset} & \textbf{Taxon} & \textbf{Novel species} &
  \textbf{Images} & \textbf{Source} \\
\midrule
ZooBirds & Aves & 412 & 24{,}720 & eBird, Macaulay \\
ZooInsects & Insecta & 638 & 31{,}900 & iNaturalist, BugGuide \\
ZooPlants & Plantae & 524 & 26{,}200 & PlantNet, herbaria \\
ZooMarine & Mixed marine & 293 & 14{,}650 & Reef Life Survey \\
\bottomrule
\end{tabular}
\end{table}

\paragraph{Citizen science field test.} We additionally evaluate on
2{,}847 field photographs submitted by 312 citizen science volunteers
through the Zoo platform, covering 187 species across all four taxa.
Images are accompanied by GPS coordinates and dates.

\subsection{Baselines}

\begin{itemize}
  \item \textbf{ProtoNet} \citep{snell2017prototypical}: Standard
    prototypical network with ResNet-50 backbone.
  \item \textbf{MAML} \citep{finn2017model}: Model-agnostic
    meta-learning with 5 inner-loop steps.
  \item \textbf{DeepEMD} \citep{zhang2020deepemd}: Earth mover's
    distance matching for few-shot.
  \item \textbf{FEAT} \citep{ye2020few}: Set-to-set function with
    transformer adaptation.
  \item \textbf{P$>$M$>$F} \citep{hu2022pushing}: Pre-train, meta-learn,
    fine-tune pipeline.
  \item \textbf{DINOv2+kNN}: DINOv2 features with $k$-nearest-neighbor
    classification (no meta-learning).
\end{itemize}

All baselines use ViT-L/14 backbone for fair comparison except where
architecturally infeasible.

\subsection{Results: Standard Few-Shot}

\begin{table}[H]
\centering
\caption{Few-shot classification accuracy (\%) on novel species. Mean
$\pm$ 95\% CI over 10{,}000 episodes.}
\label{tab:fewshot}
\begin{tabular}{l cc cc cc cc}
\toprule
& \multicolumn{2}{c}{\textbf{ZooBirds}} &
  \multicolumn{2}{c}{\textbf{ZooInsects}} &
  \multicolumn{2}{c}{\textbf{ZooPlants}} &
  \multicolumn{2}{c}{\textbf{ZooMarine}} \\
\cmidrule(lr){2-3} \cmidrule(lr){4-5} \cmidrule(lr){6-7}
  \cmidrule(lr){8-9}
\textbf{Method} & 1-shot & 5-shot & 1-shot & 5-shot &
  1-shot & 5-shot & 1-shot & 5-shot \\
\midrule
ProtoNet & 62.3 & 79.8 & 54.7 & 73.2 & 58.1 & 76.4 & 60.9 & 78.1 \\
MAML & 64.1 & 81.3 & 56.2 & 74.8 & 59.7 & 77.9 & 62.4 & 79.6 \\
DeepEMD & 67.8 & 84.1 & 59.3 & 78.6 & 63.2 & 81.2 & 65.7 & 82.3 \\
FEAT & 69.2 & 85.4 & 61.1 & 80.3 & 65.1 & 82.7 & 67.3 & 83.8 \\
P$>$M$>$F & 72.6 & 87.9 & 64.8 & 83.1 & 68.4 & 85.3 & 70.1 & 86.2 \\
DINOv2+kNN & 71.3 & 86.2 & 63.1 & 81.4 & 66.7 & 83.9 & 68.8 & 84.7 \\
\midrule
\textbf{TaxoNet} & \textbf{80.1} & \textbf{92.4} & \textbf{74.2} &
  \textbf{89.1} & \textbf{78.9} & \textbf{91.6} & \textbf{80.4} &
  \textbf{91.7} \\
\bottomrule
\end{tabular}
\end{table}

TaxoNet achieves the highest accuracy across all datasets and shot
settings. The largest improvement over the next-best method (P$>$M$>$F)
is on ZooInsects (+9.4\% 1-shot), consistent with insects having the
deepest taxonomic hierarchies and most confusable species pairs.

\subsection{Results: Citizen Science Field Test}

\begin{table}[H]
\centering
\caption{Field identification accuracy (\%) on citizen science
photographs. ``+Geo'' adds geographic and temporal priors.}
\label{tab:citizen}
\begin{tabular}{lcccc}
\toprule
\textbf{Method} & \textbf{Top-1} & \textbf{Top-3} & \textbf{Top-5} &
  \textbf{Top-1 +Geo} \\
\midrule
iNaturalist CV & 82.3 & 91.4 & 94.1 & 87.6 \\
DINOv2+kNN & 76.8 & 88.2 & 92.3 & 83.1 \\
P$>$M$>$F (5-shot) & 79.4 & 89.7 & 93.2 & 85.8 \\
\textbf{TaxoNet (5-shot)} & \textbf{84.7} & \textbf{93.1} &
  \textbf{96.2} & \textbf{94.7} \\
\bottomrule
\end{tabular}
\end{table}

With geographic priors, TaxoNet achieves 94.7\% top-1 accuracy on
citizen science data, exceeding iNaturalist's trained classifier on
overlapping species. This demonstrates the practical value of combining
few-shot visual recognition with ecological metadata.

\subsection{Ablation Studies}

\begin{table}[H]
\centering
\caption{Ablation study on ZooBirds (5-way 5-shot accuracy \%).}
\label{tab:ablation}
\begin{tabular}{lc}
\toprule
\textbf{Configuration} & \textbf{Accuracy} \\
\midrule
Full TaxoNet & \textbf{92.4} \\
\midrule
\quad -- Taxonomic prototypes (species only) & 88.7 \\
\quad -- Rank attention (uniform weights) & 90.1 \\
\quad -- Domain-adaptive pre-training (ImageNet init) & 86.3 \\
\quad -- Taxonomic episode sampling (uniform) & 89.8 \\
\quad -- Spatial attention (CLS token only) & 90.6 \\
\quad -- Taxonomic triplet loss & 91.2 \\
\bottomrule
\end{tabular}
\end{table}

Domain-adaptive pre-training provides the largest single contribution
(+6.1\%), followed by taxonomic prototypes (+3.7\%) and taxonomic
episode sampling (+2.6\%).

\subsection{Rank-Attention Analysis}

\begin{table}[H]
\centering
\caption{Learned rank-attention weights $\alpha_r$ for different query
difficulty levels.}
\label{tab:rank_attention}
\begin{tabular}{lccc}
\toprule
\textbf{Query difficulty} & $\alpha_{\text{family}}$ &
  $\alpha_{\text{genus}}$ & $\alpha_{\text{species}}$ \\
\midrule
Easy (inter-family) & 0.62 & 0.24 & 0.14 \\
Medium (inter-genus) & 0.18 & 0.49 & 0.33 \\
Hard (intra-genus) & 0.07 & 0.28 & 0.65 \\
\bottomrule
\end{tabular}
\end{table}

The rank-attention module adaptively shifts weight toward species-level
prototypes for hard (intra-genus) queries and toward family-level
prototypes for easy (inter-family) queries.


% ══════════════════════════════════════════════════════════════════════════════
\section{System Design for Citizen Science}
\label{sec:system}
% ══════════════════════════════════════════════════════════════════════════════

\subsection{Deployment Architecture}

TaxoNet is deployed as a REST API service behind a geographic routing
layer. When a citizen scientist submits a photograph:

\begin{enumerate}
  \item \textbf{Candidate filtering}: Geographic and temporal priors
    reduce the candidate species set to 50--200 plausible species.
  \item \textbf{Feature extraction}: The ViT backbone extracts visual
    features in approximately 45\,ms on an NVIDIA T4 GPU.
  \item \textbf{Few-shot matching}: Taxonomic prototypical matching
    against the candidate set runs in approximately 12\,ms.
  \item \textbf{Result presentation}: Top-5 predictions with confidence
    scores and diagnostic field marks are returned to the user.
\end{enumerate}

Total inference latency is under 100\,ms per image.

\subsection{Prototype Management}

Species prototypes are computed offline from curated reference images and
stored in a vector database indexed by taxonomy and geography. New
species can be added by computing and inserting their prototype
embeddings---no model retraining is required.

\subsection{Active Learning Loop}

Citizen science identifications that disagree with model predictions
are flagged for expert review. Confirmed identifications update
prototypes:
\begin{equation}
  \mathbf{c}_{\text{new}}^{(s)} = (1 - \eta) \mathbf{c}_{\text{old}}^{(s)}
  + \eta \cdot \phi(x_{\text{confirmed}})
\end{equation}
where $\eta = 0.01$ is the update rate.


% ══════════════════════════════════════════════════════════════════════════════
\section{Discussion}
\label{sec:discussion}
% ══════════════════════════════════════════════════════════════════════════════

\subsection{Why Taxonomy Matters}

The success of taxonomic prototypes reflects a biological reality:
species within the same genus share morphological features due to common
ancestry. By explicitly modeling this hierarchical structure, TaxoNet
leverages phylogenetically informative features even when species-level
examples are scarce. The rank-attention mechanism acts as a soft decision
tree, routing queries to the appropriate level of taxonomic resolution
based on visual evidence quality.

\subsection{Limitations}

\paragraph{Cryptic species.} Species complexes that are
morphologically indistinguishable cannot be resolved by visual features
alone. Integration with molecular data (eDNA) or acoustic features is
needed.

\paragraph{Phenological variation.} Some species change appearance
dramatically across seasons. Explicit temporal prototype modulation
would be more principled than our current update mechanism.

\paragraph{Geographic bias.} Pre-training data is biased toward North
America and Europe. Performance on tropical and Southern Hemisphere
species may be lower due to representation gaps.

\subsection{Ethical Considerations}

\begin{itemize}
  \item \textbf{Poaching risk}: We implement geographic fuzzing for
    threatened species locations.
  \item \textbf{Over-reliance}: We display confidence calibration and
    encourage expert confirmation for low-confidence predictions.
  \item \textbf{Data sovereignty}: We follow the Nagoya Protocol and
    CARE principles for indigenous data governance.
\end{itemize}


% ══════════════════════════════════════════════════════════════════════════════
\section{Implementation Details}
\label{sec:implementation}
% ══════════════════════════════════════════════════════════════════════════════

\subsection{Software and Hardware}

TaxoNet is implemented in PyTorch 2.2 with the \texttt{timm} library for
ViT models. Training was conducted on 8 NVIDIA A100 GPUs for
approximately 72 hours. Inference is deployed on NVIDIA T4 GPUs behind a
Kubernetes load balancer.

\subsection{Data Augmentation}

Training episodes apply random resized crop (scale 0.4--1.0), horizontal
flip, color jitter, random Gaussian blur, random grayscale (probability
0.1), and cutout with one $64 \times 64$ hole.

\subsection{Confidence Calibration}

We calibrate prediction confidence using temperature scaling
\citep{guo2017calibration}. Post-calibration, expected calibration error
(ECE) drops from 8.3\% to 2.1\%.


% ══════════════════════════════════════════════════════════════════════════════
\section{Future Work}
\label{sec:future}
% ══════════════════════════════════════════════════════════════════════════════

\begin{enumerate}
  \item \textbf{Multimodal fusion}: Combining visual features with audio
    and eDNA sequence embeddings for visually cryptic species.
  \item \textbf{Continual meta-learning}: Adapting as new species are
    described without catastrophic forgetting.
  \item \textbf{Federated training}: Training on distributed citizen
    science data without centralizing raw images, governed by Zoo DAO.
  \item \textbf{Zero-shot identification}: Leveraging taxonomic
    descriptions and field guide text via vision-language models.
  \item \textbf{Decentralized inference}: Deploying TaxoNet on edge
    devices and the Zoo decentralized compute network.
\end{enumerate}


% ══════════════════════════════════════════════════════════════════════════════
\section{Conclusion}
\label{sec:conclusion}
% ══════════════════════════════════════════════════════════════════════════════

We have presented TaxoNet, a taxonomy-aware few-shot visual recognition
system for biodiversity monitoring. By incorporating biological
classification structure into meta-learning through taxonomic prototypical
networks and rank-attention fusion, TaxoNet achieves state-of-the-art
few-shot accuracy across four major taxonomic groups. Integration with
geographic and phenological priors further boosts practical field
identification to 94.7\% accuracy. Through Zoo Labs Foundation's
commitment to open science, TaxoNet's models, benchmarks, and deployment
infrastructure are freely available to the conservation community.


% ══════════════════════════════════════════════════════════════════════════════
% References
% ══════════════════════════════════════════════════════════════════════════════

\bibliographystyle{plainnat}

\begin{thebibliography}{33}

\bibitem[Akata et~al.(2015)]{akata2015label}
Akata, Z., Reed, S., Walter, D., Lee, H., and Schiele, B.
\newblock Evaluation of output embeddings for fine-grained image classification.
\newblock \emph{CVPR}, pp. 2927--2936, 2015.

\bibitem[Beery et~al.(2021)]{beery2021iwildcam}
Beery, S., Cole, E., Parker, J., Perona, P., and Winner, K.
\newblock Species classification in camera traps: The iWildCam 2021 challenge.
\newblock \emph{FGVC Workshop at CVPR}, 2021.

\bibitem[Bengio et~al.(2009)]{bengio2009curriculum}
Bengio, Y., Louradour, J., Collobert, R., and Weston, J.
\newblock Curriculum learning.
\newblock \emph{ICML}, pp. 41--48, 2009.

\bibitem[Brown et~al.(2020)]{brown2020language}
Brown, T., Mann, B., Ryder, N., et~al.
\newblock Language models are few-shot learners.
\newblock \emph{NeurIPS}, 33:1877--1901, 2020.

\bibitem[Chang et~al.(2021)]{chang2021your}
Chang, D., Pang, K., Zheng, Y., et~al.
\newblock Your ``flamingo'' is my ``bird'': Fine-grained, or not.
\newblock \emph{CVPR}, pp. 11476--11485, 2021.

\bibitem[Christin et~al.(2019)]{christin2019applications}
Christin, S., Hervet, {\'E}., and Lecomte, N.
\newblock Applications for deep learning in ecology.
\newblock \emph{Methods in Ecology and Evolution}, 10(10):1632--1644, 2019.

\bibitem[Costello et~al.(2013)]{costello2013can}
Costello, M.J., May, R.M., and Stork, N.E.
\newblock Can we name Earth's species before they go extinct?
\newblock \emph{Science}, 339(6118):413--416, 2013.

\bibitem[Deng et~al.(2014)]{deng2014large}
Deng, J., Ding, N., Jia, Y., et~al.
\newblock Large-scale object classification using label relation graphs.
\newblock \emph{ECCV}, pp. 48--64, 2014.

\bibitem[Dosovitskiy et~al.(2021)]{dosovitskiy2021image}
Dosovitskiy, A., Beyer, L., Kolesnikov, A., et~al.
\newblock An image is worth 16x16 words: Transformers for image recognition.
\newblock \emph{ICLR}, 2021.

\bibitem[Finn et~al.(2017)]{finn2017model}
Finn, C., Abbeel, P., and Levine, S.
\newblock Model-agnostic meta-learning for fast adaptation of deep networks.
\newblock \emph{ICML}, pp. 1126--1135, 2017.

\bibitem[Go{\"e}au et~al.(2022)]{goeau2022overview}
Go{\"e}au, H., Bonnet, P., and Joly, A.
\newblock Overview of PlantCLEF 2022: Image-based plant identification.
\newblock \emph{CLEF}, 2022.

\bibitem[Guo et~al.(2017)]{guo2017calibration}
Guo, C., Pleiss, G., Sun, Y., and Weinberger, K.Q.
\newblock On calibration of modern neural networks.
\newblock \emph{ICML}, pp. 1321--1330, 2017.

\bibitem[He et~al.(2022)]{he2022transfg}
He, J., Chen, J.N., Liu, S., et~al.
\newblock TransFG: A transformer architecture for fine-grained recognition.
\newblock \emph{AAAI}, pp. 852--860, 2022.

\bibitem[Hu et~al.(2018)]{hu2018squeeze}
Hu, J., Shen, L., and Sun, G.
\newblock Squeeze-and-excitation networks.
\newblock \emph{CVPR}, pp. 7132--7141, 2018.

\bibitem[Hu et~al.(2022)]{hu2022pushing}
Hu, S.X., Li, D., St{\"u}hmer, J., Kim, M., and Hospedales, T.M.
\newblock Pushing the limits of simple pipelines for few-shot learning.
\newblock \emph{CVPR}, pp. 9068--9077, 2022.

\bibitem[Kahl et~al.(2021)]{kahl2021birdnet}
Kahl, S., Wood, C.M., Eibl, M., and Klinck, H.
\newblock BirdNET: A deep learning solution for avian diversity monitoring.
\newblock \emph{Ecological Informatics}, 61:101236, 2021.

\bibitem[Li et~al.(2019)]{li2019large}
Li, A., Luo, T., Lu, Z., Xiang, T., and Wang, L.
\newblock Large-scale few-shot learning: Knowledge transfer with class hierarchy.
\newblock \emph{CVPR}, pp. 7212--7220, 2019.

\bibitem[Lin et~al.(2015)]{lin2015bilinear}
Lin, T.Y., RoyChowdhury, A., and Maji, S.
\newblock Bilinear CNN models for fine-grained visual recognition.
\newblock \emph{ICCV}, pp. 1449--1457, 2015.

\bibitem[Liu et~al.(2020)]{liu2020many}
Liu, B., Kang, H., Li, H., Hua, G., and Vasconcelos, N.
\newblock Few-shot open-set recognition using meta-learning.
\newblock \emph{CVPR}, pp. 8798--8807, 2020.

\bibitem[Lorieul et~al.(2022)]{lorieul2022overview}
Lorieul, T., Cole, E., Deneu, B., et~al.
\newblock Overview of GeoLifeCLEF 2022: Predicting species presence.
\newblock \emph{CLEF}, 2022.

\bibitem[Norouzzadeh et~al.(2018)]{norouzzadeh2018automatically}
Norouzzadeh, M.S., Nguyen, A., Kosmala, M., et~al.
\newblock Automatically identifying wild animals in camera-trap images.
\newblock \emph{PNAS}, 115(25):E5716--E5725, 2018.

\bibitem[Oquab et~al.(2024)]{oquab2024dinov2}
Oquab, M., Darcet, T., Moutakanni, T., et~al.
\newblock DINOv2: Learning robust visual features without supervision.
\newblock \emph{TMLR}, 2024.

\bibitem[Snell et~al.(2017)]{snell2017prototypical}
Snell, J., Swersky, K., and Zemel, R.
\newblock Prototypical networks for few-shot learning.
\newblock \emph{NeurIPS}, pp. 4077--4087, 2017.

\bibitem[Steenweg et~al.(2017)]{steenweg2017scaling}
Steenweg, R., Hebblewhite, M., Kays, R., et~al.
\newblock Scaling-up camera traps: monitoring biodiversity with sensor networks.
\newblock \emph{Frontiers in Ecology and the Environment}, 15(1):26--34, 2017.

\bibitem[Sung et~al.(2018)]{sung2018learning}
Sung, F., Yang, Y., Zhang, L., et~al.
\newblock Learning to compare: Relation network for few-shot learning.
\newblock \emph{CVPR}, pp. 1199--1208, 2018.

\bibitem[Van~Horn et~al.(2015)]{van2015building}
Van~Horn, G., Branson, S., Farrell, R., et~al.
\newblock Building a bird recognition app with citizen scientists.
\newblock \emph{CVPR}, pp. 595--604, 2015.

\bibitem[Van~Horn et~al.(2021)]{van2021benchmarking}
Van~Horn, G., Cole, E., Beery, S., et~al.
\newblock Benchmarking representation learning for natural world images.
\newblock \emph{CVPR}, pp. 12884--12893, 2021.

\bibitem[Vinyals et~al.(2016)]{vinyals2016matching}
Vinyals, O., Blundell, C., Lillicrap, T., and Welling, M.
\newblock Matching networks for one shot learning.
\newblock \emph{NeurIPS}, pp. 3630--3638, 2016.

\bibitem[Wang et~al.(2020)]{wang2020generalizing}
Wang, Y., Yao, Q., Kwok, J.T., and Ni, L.M.
\newblock Generalizing from a few examples: A survey on few-shot learning.
\newblock \emph{ACM Computing Surveys}, 53(3):1--34, 2020.

\bibitem[Wei et~al.(2019)]{wei2019piecewise}
Wei, X.S., Wang, P., Liu, L., Shen, C., and Wu, J.
\newblock Piecewise classifier mappings for few-shot fine-grained recognition.
\newblock \emph{IEEE TIP}, 28(12):6116--6125, 2019.

\bibitem[Ye et~al.(2020)]{ye2020few}
Ye, H.J., Hu, H., Zhan, D.C., and Sha, F.
\newblock Few-shot learning via saliency-guided hallucination of samples.
\newblock \emph{CVPR}, pp. 770--779, 2020.

\bibitem[Zhang et~al.(2020)]{zhang2020deepemd}
Zhang, C., Cai, Y., Lin, G., and Shen, C.
\newblock DeepEMD: Few-shot image classification with differentiable EMD.
\newblock \emph{CVPR}, pp. 12203--12213, 2020.

\bibitem[Zheng et~al.(2017)]{zheng2017learning}
Zheng, H., Fu, J., Mei, T., and Luo, J.
\newblock Learning multi-attention CNN for fine-grained image recognition.
\newblock \emph{ICCV}, pp. 5209--5217, 2017.

\end{thebibliography}

% ══════════════════════════════════════════════════════════════════════════════
\appendix
\section{Supplementary: Per-Taxon Analysis}
\label{app:per_taxon}
% ══════════════════════════════════════════════════════════════════════════════

\begin{table}[H]
\centering
\caption{Detailed per-taxon performance breakdown (5-way 5-shot \%).}
\label{tab:per_taxon}
\begin{tabular}{llcc}
\toprule
\textbf{Dataset} & \textbf{Difficulty} & \textbf{TaxoNet} &
  \textbf{P$>$M$>$F} \\
\midrule
\multirow{3}{*}{ZooBirds}
  & Inter-family & 97.8 & 96.1 \\
  & Inter-genus & 93.6 & 88.4 \\
  & Intra-genus & 86.2 & 78.9 \\
\midrule
\multirow{3}{*}{ZooInsects}
  & Inter-family & 96.4 & 94.2 \\
  & Inter-genus & 90.1 & 83.7 \\
  & Intra-genus & 81.3 & 71.6 \\
\midrule
\multirow{3}{*}{ZooPlants}
  & Inter-family & 97.1 & 95.3 \\
  & Inter-genus & 92.8 & 86.1 \\
  & Intra-genus & 84.7 & 74.8 \\
\midrule
\multirow{3}{*}{ZooMarine}
  & Inter-family & 97.3 & 95.7 \\
  & Inter-genus & 92.4 & 87.0 \\
  & Intra-genus & 85.1 & 76.2 \\
\bottomrule
\end{tabular}
\end{table}

TaxoNet's advantage is most pronounced for intra-genus discrimination
(+7.3\% to +9.9\% over P$>$M$>$F), precisely the setting where
taxonomic prototypes provide the most value.

\section{Supplementary: Prototype Visualization}
\label{app:visualization}

Analysis of learned prototypes using t-SNE reveals clear taxonomic
clustering: species prototypes cluster tightly within genera, genera form
distinct sub-clusters within families, and families separate cleanly in
the embedding space. The taxonomic prototypical network naturally
organizes the embedding space to reflect biological relationships,
emerging from the hierarchical episode sampling and multi-rank loss.

Spatial attention maps show that TaxoNet focuses on taxonomically
informative features: wing bar patterns for warbler identification,
venation patterns for butterfly classification, leaf margin shapes for
plant identification, and operculum structure for marine gastropod
classification. These attention patterns are consistent with the
diagnostic characters used by human taxonomists.

\end{document}
